\documentclass{article}
\usepackage{bm}
%Para fotos y eso
\usepackage{graphicx}
%Para listas
\usepackage{enumitem}
%Para fórmulas y símbolos
\usepackage{amsmath}
\usepackage{amsfonts}
\usepackage{amssymb}
\usepackage{aliascnt}
%Para teoremas y demostraciones
\usepackage{amsthm}
%Para cuadraditos :)
\usepackage{mdframed}
%Headers y footers
\usepackage{fancyhdr}
%Para manejar el formato del papel
\usepackage{geometry}
%Manejo de color
\usepackage{color}
%Hipervínculos
\usepackage[colorlinks,linkcolor=blue]{hyperref}
\usepackage[nameinlink]{cleveref}
%cajitas
\usepackage[most]{tcolorbox}
\tcbuselibrary{listings, breakable, skins, theorems}
\usepackage{float}
%tikz
\usepackage{tikz,tikz-cd}

\usepackage{titlesec}
\titleformat{\chapter}[hang]{\Huge\bfseries}{Capítulo \thechapter}{1em}{}

\usepackage{graphicx,xcolor}

\usetikzlibrary{arrows.meta, positioning,calc, shapes.geometric}
%Hace setup de todos los hipervínculos del proyecto
\hypersetup{
    colorlinks=true,
    linktoc=all,
    linkcolor=black,
}
\geometry{letterpaper,margin=1in}
%SÍMBOLOS MATEMÁTICOS
%Conjuntos N,Z,Q,R,C:
\newcommand{\N}{\mathbb{N}}
    \newcommand{\Z}{\mathbb{Z}}
    \newcommand{\Q}{\mathbb{Q}}
    \newcommand{\R}{\mathbb{R}}
    \newcommand{\C}{\mathbb{C}}
    \newcommand{\F}{\mathbb{F}}
    \newcommand{\E}{\mathbb{E}}
    \newcommand{\Pp}{\mathbb{P}}
%Mathcal
\newcommand{\gL}{\mathcal{L}}
\newcommand{\gF}{\mathcal{F}}
\newcommand{\gU}{\mathcal{U}}
\newcommand{\gQ}{\mathcal{Q}}
\newcommand{\gG}{\mathcal{G}}
\newcommand{\gC}{\mathcal{C}}
\newcommand{\gT}{\mathcal{T}}
\newcommand{\gP}{\mathcal{P}}
\newcommand{\gB}{\mathcal{B}}
\newcommand{\gS}{\mathcal{S}}
\newcommand{\gO}{\mathcal{O}}
%Conjunto potencia
\newcommand{\powerset}[1]{\mathcal{P}(#1)}
%Para conjuntos
\newcommand{\set}[1]{\left\{#1\right\}}
%Función f:A->B
\newcommand{\func}[3]{#1: #2 \to #3}
%Vector cero
\newcommand{\zerov}{\overline{0}}
%Una lista indexada
\newcommand{\vectorlist}[2]{#1_1,\dots,#1_{#2}}
%Una combinación lineal
\newcommand{\lcombination}[3]{#1_1#2_1+#1_2#2_2+\dots+#1_{#3}#2_{#3}}

%Una lgualdad indexada de elementos, el tercer elemento es distinto del resto
\newcommand{\equallist}[3]{#1_1=#1_2=\dots=#1_{#2}=#3}
%Vector en F^2
\newcommand{\twovector}[2]{\begin{pmatrix}
#1\\
#2
\end{pmatrix}}
%Vector columna con n entradas
\newcommand{\coordvector}[2]{\begin{pmatrix}
#1_1\\
#1_2\\
\vdots\\
#1_{#2}
\end{pmatrix}}
%Inverso
\newcommand{\minv}[1]{#1^{-1}}
%Suma
\newcommand{\summ}[3]{\sum_{#1=#2}^{#3}}
%Representación matricial de una transformación lineal
\newcommand{\matrep}[3]{[#1]_{#2}^{#3}}
%Subsección no enumerada se añade a la tabla de contenidos.
\newcommand{\nsubsection}[1]{
\subsection*{#1}
\addcontentsline{toc}{subsection}{#1}
}

% Enteros módulo n
\newcommand{\Zn}[1]{\Z/#1\Z}
%un m-ciclo
\newcommand{\mcycle}[2]{(#1_1\spa #1_2\spa \dots \spa #1_{#2})}
%grupo cíclico
\newcommand{\cyclicgrp}[1]{\langle #1 \rangle}

\newcommand{\charsbgrp}{\ \text{char}\ }
\usepackage{graphicx} % Required for inserting images

\title{Proyecto Decimo Problema}
\author{Juan Camilo Solano, Nicolas Hernandez Derchhhhhh, Samuel Montoya Salazar}
\date{March 2026}
% Define English "lemma" environment to match usage \begin{lemma}...
\newtheorem{lemma}{Lema}
\newenvironment{dem}{\par\noindent\textbf{Dem.}\par}{\hfill$\square$\par}
\newenvironment{teo}{\par\noindent\textbf{Teorema.}\par}{\hfill$\square$\par}
\newenvironment{lem}{\par\noindent\textbf{Lema.}\par}{\hfill$\square$\par}
\newenvironment{deff}{\par\noindent\textbf{Definición.}\par}{\hfill$\square$\par}

\begin{document}
\maketitle

\begin{abstract}
    En este documento buscamos hacer una demostracion a cabalidad de la existencia de un algoritmo polinomial
    deterministico para determinar si un numero es primo o no.
\end{abstract}

\section{Historia de como encontrar primos}
A lo largo de la historia, los matematicos han desarrollado un
interes intrinseco por encontrar numeros primos. Sus propiedades
abarcan todas las areas de las matematicas, llegando desde su
origen la teoria de numeros hasta a areas en el algebra y 
geometria. Es por esto que desde los griegos, se busca una manera
de encontrar primos de diferentes maneras.\\
Con el tiempo los matematicos encontraron muchos metodos. El mas antiguo
y mas sencillo de aplicar es el metodo de la criba de Eratostenes, que
consiste en calcular la raiz cuadrada del numero del que se quiere
si es primo o no. Si la raiz es un entero entonces el numero era un cuadrado
y por tanto no era primo, de lo contrario se calcula el piso de la raiz y 
miran los numeros primos menores a este numero, si alguno de estos numeros 
divide al numero quiere decir que el numero no es primo.\\
Pero hay algo que todos estos metodos pasan por alto,
y es que el numero de numeros primos menores a un numero n es aproximadamente 
$\frac{\sqrt{n}}{\log(\sqrt{n})}$, lo que hace que el metodo de la criba de 
Eratostenes tenga una complejidad de $\it{O}(\frac{\sqrt{n}}{\log(\sqrt{n})})$ lo 
que no es de complejidad polinomial.\\
Lo que quiere decir esto es que, aunque el metodo de un resultado
correcto, no es un metodo eficiente para encontrar numeros primos.
Otro metodo que se usaba antes era el metodo usando el pequeño teorema 
de Fermat, que dice que si p es un numero primo y a es un numero entero 
que no es divisible por p, entonces $a^{p-1} \equiv 1 \mod p$. Este metodo
se creia que servia para encontrar primos pues el algoritmo sera
elegir un numero a y calcular $a^{n-1} \mod n$, si el resultado es diferente de 1, entonces n no es primo.
y se prueba con los a menores a n. Sin embargo, este metodo no es correcto pero si polinomial a diferencia del anterior metodo, 
pues existen numeros compuestos llamados pseudoprimos de Fermat o los numeros de Carmichael que cumplen con 
esta propiedad para todos los valores de a con los que sean coprimos con n.\\
Si bien existen muchos criterios para encontrar numeros primos, como el test de  Willson, la idea principal del algoritmo
es que el criterio en el cual se basa el algoritmo tiene una cota lo suficientemente buena para
lograr una complejidad polinomial menor a la que ya se tenia.

\section{Primes is in P}

Para iniciar la demostracion de este algoritmo, vamos a demostrar el criterio esencial el cual
es muy similar al Pequeño teorema de Fermat, pero con una cota mucho mas estricta.\\
\begin{teo}
    Sea $a\in \Z, n\in \N, n\geq 2$ y $(a,n)=1$, entonces n es primo si y solo si $(x+a)^n \equiv x^n + a \mod n(x)$
\end{teo}

\begin{dem}
    Para demostrar la implicacion directa suponga que n es primo entonces por teorema binomial se tiene que:
\begin{align*}
(x+a)^n &= \sum_{k=0}^n \binom{n}{k} x^{n-k} a^k \\
&= x^n + a^n \mod n(x)
\end{align*}
pues cada termino de $\binom{n}{k}$ con $1\leq k \leq n-1$ es divisible por n pues siempre esta el coeficiente de n en el numerador.
Esto deja que $(x+a)^n \equiv x^n + a^n \mod n(x)$ y por pequeño teorema de Fermat se tiene que $a^n \equiv a \mod n$ y por tanto $(x+a)^n \equiv x^n + a \mod n(x)$.\\
Para la suficiencia suponga que n es compuesto, entonces existe un primo $q$ tal que $q|n$ y considere $k\in \N$ tal que $q^k||n$. Note ahora que $q^k \not| \binom{n}{q}$ pues 
$\binom{n}{q}=\frac{n!}{q!(n-q!)}=\frac{(n-1)!m}{(q-1)!(n-q)!}$ dejando que $q^{k-1}|m \to q^k\not|\binom{n}{q}$. Como $(a,n)=1$ entonces $q^k\not|a \to q^k\not|a^r \forall r\in \N$, 
dejando que como $(q^k,a^{n-q})=1$ entonces el coeficiente en la expancion binomial del termino $x^q$ no va a ser 0 modulo n y por tanto $(x+a)^n - x^n + a a \neq 0 \mod n(x)$ dando a lugar a una contradiccion

\end{dem}

\subsection{Notaciones}
Defina el orden multiplicativo de n modulo r como el menor entero positivo k tal que $n^k \equiv 1 \mod r$ equivalente a el orden en $\Z_r^*$ y se denota por $O_r(n)$, 
es importante aclarar que el orden solo tiene sentido si $(n,r)=1$.\\
$\Z_n$ denota el anillo con elementos modulo n, y $\F_p$ denota el campo con p elementos, es decir $\Z_p$ con p primo.\\
Si $h(x)$ es un polinomio de grado d entonces $\F[x]/(h(x))$ denota el cuerpo finito de orden $p^d$. Decimos que $f(x) \equiv g(x) \mod h(x), n$ para representar que $f(x)$ y $g(x)$ son equivalentes en el anillo $\Z_n[x]/(h(x))$.\\
$\phi(n)$ denota la funcion de Euler totient, es decir el numero de enteros positivos menores a n que son coprimos con n.\\


Ahora con este resultado podemos empezar a entrar mas a profundidad en el algoritmo. Si bien es un resultado muy importante en las matematicas
su importancia no es proporcional con su dificultad a diferencia de muchas otras cosas dentro del area. Es un poco contraintuitivo que 
un resultado asi se demorase tanto en ser descubierto puesto que el conocimiento de aqui en adelante para demostrar el algoritmo es un conocimiento
de pregrado de matematicas y posiblemente menos para ser demostrado en su totalidad. A tal efecto que quienes lo demostraron fueron estudiantes.\\


\subsection{El Algoritomo}

\begin{center}
\begin{tabular}{|p{0.85\textwidth}|}
\hline
\\
\texttt{Entrada:}\quad entero $n > 1$.\\[6pt]
\texttt{1.}\quad Si $(n = a^b$ Para $a \in \N$ y $b > 1)$, retorne \texttt{Compuesto}.\\[6pt]
\texttt{2.}\quad Encuentre el $r$ mas pequeño tal que $\mathrm{o}_r(n) > \log^2 n$.\\[6pt]
\texttt{3.}\quad Si $1 < (a, n) < n$ Para algun $a \leq r$, retorne \texttt{Compuesto}.\\[6pt]
\texttt{4.}\quad Si $n \leq r$, retorne \texttt{Primo}.$^1$\\[6pt]
\texttt{5.}\quad Desde $a = 1$ hasta $\lfloor\sqrt{\phi(r)}\log n\rfloor$ haga lo siguiente:\\[4pt]
\qquad\qquad Si $\left((X + a)^n \neq X^n + a \pmod{X^r - 1,\, n}\right)$, retorne \texttt{Compuesto};\\[6pt]
\texttt{6.}\quad retorne \texttt{Primo}.\\[6pt]
\hline
\end{tabular}
\end{center}

Ahora vamos a empezar a demostrar los resultados necesarios que nos llevaran al resultado final.\\

\subsection{Demostracion del Algoritmo}
\begin{lem}
Sea $L(m) = \text{mcm}(1, 2, \dots, m)$. Para todo $m \ge 7$, se cumple que:
\[ L(m) \ge 2^m \]
\end{lem}

\begin{dem}
La demostración se basa en el estudio de la integral:
$ I_n = \int_{0}^{1} x^n (1-x)^n dx $
Es un resultado conocido que $L(2n+1) I_n$ es un entero positivo. Dado que $I_n = \frac{(n!)^2}{(2n+1)!}$, tenemos:
\begin{equation}
L(2n+1) \ge \frac{(2n+1)!}{(n!)^2} = (2n+1) \binom{2n}{n}
\end{equation}
Utilizando la desigualdad para el coeficiente binomial central $\binom{2n}{n} \ge \frac{4^n}{2n}$, para $n \ge 1$, se obtiene:
\[
L(2n+1) \ge (2n+1)\frac{4^n}{2n} = \left(1+\frac{1}{2n}\right)2^{2n}.
\]
Para $m = 2n+1$, esto implica $L(m) > 2^{m-1}$. Un refinamiento de este método mediante el uso de integrales de la forma $\int_{0}^{1} x^{n+k}(1-x)^n dx$ permite fortalecer la cota. 

Específicamente, para $m \ge 7$, procedemos por inducción o verificación directa. 
Para el caso base $m=7$:
\[ L(7) = 420 > 2^7 = 128 \]
Para el paso inductivo, se utiliza el hecho de que $L(m+1) = L(m) \cdot \frac{m+1}{\text{mcd}(L(m), m+1)}$. El crecimiento de la función $\psi(m) = \ln L(m)$ está acotado inferiormente por $m \ln 2$ para $m \ge 7$ según el análisis de las integrales de Chebyshev, lo que concluye que:
\[ L(m) \ge 2^m \]
\end{dem} 

Ahora vamos a demostrar las dos implicaciones del teoremaque van a terminar dando el resultado final. La primera de estas es mucho mas facil que la reciproca y por ello se siguen los siguientes resultados.\\

\begin{teo}
    Si n es primo entonces el algoritmo retorna \texttt{Primo}.
\end{teo}
\begin{dem}
    Como n es primo entonces no se cumple ni el paso 1 ni el 3 por definicion de ser primo. Ahora por el teorema 1 se tiene que n no puede cumplir el paso 5, pues si n es primo entonces se cumple que $(x+a)^n \equiv x^n + a \mod n(x)$ para todo a, y por tanto el algoritmo retorna \texttt{Primo} sin importar en que paso lo retorne.
\end{dem}
De ahora en adelante el logaritmo siempre se va a tomar en base 2 a menos que se indique lo contrario
\begin{lem}
    $\exists r \leq max\{3, \lceil \log^5 n \rceil\}$ tal que $O_r(n) > \log^2 n$.
\end{lem}
Este resultado nos es util pues para el paso 2 del algoritmo necesitamos encontrar este r y es importante que este r no sea tan grande para que el algoritmo siga siendo polinomial.\\
\begin{dem}
    Note que para $n=2,r=3$ se cumple la condicion trivialmente pues $max\{3,\lceil \log^5 n \rceil =1\} =3$ y $O_3(2)=2>1=\log^2(2)$. Suponga que $n>2$ entonces como $\log n \geq \log 3 \geq \sqrt[5]{10}$
    Entonces  $\log^5 n \geq 10$ y $\lceil\log^5 n \rceil\ge 10$. Considere $r_1,r_2,\dots,r_t$ tales que $O_{r_i}(n)\leq \log^2 n$ o $r_i|n$, entonces todos los $r_i$ dividen al siguiente producto: \\
    \begin{align}
        n\prod_{i=1}^{\lfloor\log^2n\rfloor} (n^i -1)
    \end{align}
    puesto que si $r_i|n$ se tiene trivialmente, y si $O_{r_i}(n)\leq \log^2 n$ entonces el orden de $r_i$ debe estar en el intervalo de $[1,\log^2n]$ entonces, en particular, estan todos los ordenes de todos los $r_i$
    que no dividen a n. Por otro lado, 
    \begin{align}
        P=n\prod_{i=1}^{\lfloor\log^2n\rfloor} (n^i -1)\le n^{\log^4n}\leq 2^{\log^5n}
    \end{align}
    Dado que $n^i \leq n^{\lfloor\log^2n\rfloor}$ y 
    \begin{align}
    n\prod_{i=1}^{\lfloor\log^2n\rfloor} (n^i -1)\le n\prod_{i=1}^{\lfloor\log^2n\rfloor} (n^i)\le n^{2\lfloor\log^2n\rfloor} \le n^{2\log^2n} = n^{\log^4n}
    \end{align}
    ahora por el lema 1 se tiene que $L(\lceil\log^5 n\rceil) \geq 2^{\lceil\log^5 n\rceil} > 2^{\log^5 n} = n^{\log^4 n}$, lo que implica que $n\prod_{i=1}^{\lfloor\log^2n\rfloor} (n^i -1) < L(\lceil\log^5 n\rceil)$,
    lo que implica que debe existir un $s\leq \lceil\log^5 n\rceil$ tal que $s$ no es un $r_i$, es decir, $O_s(n) > \log^2 n$ y $s$ no divide a n. Ahora si $(s,n)=1$ queda demostrado el lema,
    si $(s,n)>1$ entonces $s|n$ o $O_s(n) \leq \log^2 n$. Defina $r=\frac{s}{(s,n)}$
    veamos que $r\notin\{r_1,r_2,\dots,r_t\}$, pues si $r=r_i$ entonces hay dos casos, $r|P$ o $O_{r_i}(n)\leq \log^2 n$. 
    Si $O_{r_i}(n)\leq \log^2 n$ y $(r_i,n)=1$ entonces defina $d=(s,n)$, entonces $(r,d)=1$ pues tanto r como d dividen a P entonces como tienen 
    minimo comun divisor 1 entonces $rd|P$ lo que es una contradiccion pues $s=r(s,n)|P$ lo que es una contradiccion pues $s$ no es un $r_i$.
    Ahora si $r|n$ entonces veamos que $(r,d)=1$, suponga que no entonces existe un primo $q$ tal que $q|r$ y $q|d$, pero como $q|s,q|n$ entonces $qd|s$ y $qd|n$
    dejando asi que $qd|d$ entonces $q=1$ contradiccion entonces como $(r,d)=1$ entonces por el argumento de antes queda que $s|n$ contradiccion.\\
\end{dem}
Con este lema podemos ir avanzando en los pasos del algoritmo, sabemos que el paso 1 ya no se cumple pues de lo contrario claramente 
habria retornado compuesto. Supongamos que el paso 1 no dio como resultado compuesto, ahora por el anterior lema podemos ejecutar el 
paso 2 sin problemas pues sabemos que existe un r. Ahora como sabemos que $O_r(n)>1$ entonces por Teorema de Lagrange debe existir un 
divisor primo de n, p, tal que $O_r(p)>1$. Ahora $p>r$ pues si $p\leq r$ entonces el paso 3 del algoritmo hubiera retornado compuesto,
por otro lado, si $n\leq p$ entonces esto deberia retornar primo pues tenemos un numero cuyo orden no es potencia de ninguno de los anteriores
en un $\Z_r^*$ que es un grupo donde $n$ tiene orden $O_r(n)$ talque $O_r(p)|O_r(n)$ pero $O_r(p)$ es el orden del gurpo entonces n 
tiene que ser primo obligatoriamente.\\
Ahora tenemos que $p\ge r$ y por el paso 2 que $(n,r)=1$, de ahora en adelante tanto p como r van a estar fijos.
Sea $\it{l}=\lfloor \sqrt{\phi(r)}\log n \rfloor$, el algoritmo ahora en el paso 5 va a verificar $\it l$ ecuaciones.
Suponga que el algoritmo no retorna compuesto en el paso 5 entonces tenemos que 

\begin{align}
    (x+a)^n \equiv x^n + a \mod(x^r - 1, n), \forall a \in[0,\it{l}]
\end{align}

Pero como $p|n$ entonces tenemos que 

\begin{align}
    (x+a)^n \equiv x^n + a \mod(x^r - 1, p), \forall a \in[0,\it{l}]
\end{align}

Ahora por el teorema 1 se tiene que 

\begin{align}
    (x+a)^p \equiv x^p + a \mod(x^r - 1, p), \forall a \in[0,\it{l}]
\end{align}

Como en un cuerpo de caracterisitica p, elevar a la p es inyectivo, pues
\begin{align}
    f(x) = g(x) \to f(x) - g(x) = 0 \to (f(x) - g(x))^p = 0 \to f(x)^p - g(x)^p = 0  \to f(x)^p = g(x)^p \mod p
\end{align}

Entonces por las ecuaciones anteriores note que 

\begin{align}
    (x+a)^{(n/p)p} = (x+a)^n = x^n + a = x^{(n/p)p} + a = (x^{n/p}+a)^p\mod(x^r - 1, p), \forall a \in[0,\it{l}] 
\end{align}
Dejando que 
\begin{align}
    (x+a)^{(n/p)p} \equiv x^{(n/p)p} + a \mod(x^r - 1, p), \forall a \in[0,\it{l}]
\end{align}
Y como sabemos que elevar a la p es inyectivo queda el $(x+a)^{(n/p)} \equiv x^{(n/p)} + a \mod(x^r - 1, p), \forall a \in[0,\it{l}]$\\
Esta propiedad la vamos a llamar introspectividad y va a ser importante mas adelante.\\

\begin{deff}
    Para un polinomio $f(x)$ y un numero $m\in \N$, se dice que $f(x)$ es introspectivo para $m$ si
    $f(x)^m \equiv f(x^m) \pmod{x^r - 1, p}$.
\end{deff}

\begin{lem}
    Si $m,n$ son enteros introspectivos para $m$ entonces $mn$ es introspectivo para $f(x)$.
\end{lem}

\begin{dem}
    Note que $f(x)^{mn} = (f(x)^m)^n \equiv f(x^m)^n  \mod(x^r - 1, p)$, remplazando $x^m=y$ tenemos que $f(y)^n \equiv f(y^n) \mod(x^r - 1, p)$, 
    remplazando $y=x^m$ tenemos que $f(x^m)^n \equiv f((x^m)^n) \mod(x^r - 1, p)$, lo que implica que $f(x)^{mn} \equiv f(x^{mn}) \mod(x^r - 1, p)$ y 
    por tanto $mn$ es introspectivo para $f(x)$.    
\end{dem}

\section{Implicaciones en la computación}

\section{Conclusiones}

\section{Referencias}


\end{document}
